% A proof of Graffiti.pc Conjecture 142 (WOWII):
% t(G) >= (2/3)g(G) + ecc(G,Periphery(G)).
\documentclass[11pt]{article}
\ifdefined\XeTeXversion\else\pdfoutput=1\fi
\usepackage[margin=1.05in]{geometry}
\usepackage{amsmath,amssymb,amsthm}
\usepackage[colorlinks=true,linkcolor=blue,citecolor=blue,urlcolor=blue]{hyperref}
\Urlmuskip=0mu plus 1mu\relax
\setlength{\emergencystretch}{2em}
\numberwithin{equation}{section}

\newtheorem{theorem}{Theorem}[section]
\newtheorem{lemma}[theorem]{Lemma}
\newtheorem{corollary}[theorem]{Corollary}
\newtheorem{proposition}[theorem]{Proposition}
\theoremstyle{remark}
\newtheorem{remark}[theorem]{Remark}

\newcommand{\tr}{\operatorname{t}}
\newcommand{\gr}{\operatorname{g}}
\newcommand{\diam}{\operatorname{diam}}
\newcommand{\ecc}{\operatorname{ecc}}
\newcommand{\Per}{\operatorname{Per}}
\newcommand{\dist}{\operatorname{dist}}

\title{Largest induced trees, girth, and distance from the periphery:\\
a proof of Graffiti.pc's Conjecture 142}
\author{Alper Ferudun\\\texttt{alper@mercurycodelab.com}}
\date{July 2026}

\begin{document}
\maketitle

\begin{abstract}
For a finite simple connected graph $G$, let $\tr(G)$ be the largest order of
an induced subgraph that is a tree, let $\gr(G)$ be the girth, let
$\Per(G)$ be the set of peripheral vertices, and put
$f(G)=\max_{x\in V(G)}d(x,\Per(G))$.  We prove
\[
  \tr(G)\ge \frac{2}{3}\gr(G)+f(G),
\]
which is Conjecture~142 of the conjecture-making program Graffiti.pc
(DeLaVi\~na, \emph{Written on the Wall II}, 2005).  In the cyclic case we
prove the stronger integral form
$\tr(G)\ge f(G)+\lceil2\gr(G)/3\rceil$.
The main ingredient is a rooted shortest-cycle lemma: in a graph of girth at
least five, any three prescribed vertices can be placed in a tree obtained
from a shortest cycle by deleting one cycle vertex and attaching an
admissible induced forest.  A three-point metric estimate then reduces the
result to exact arithmetic, with one extremal family handled separately.
\end{abstract}

\section{Introduction}

Graffiti.pc, developed by DeLaVi\~na~\cite{DeL01,DeLhistory} as a successor
of Fajtlowicz's Graffiti program~\cite{Faj88}, generates conjectural
inequalities among graph invariants.  Its conjectures are collected in
\emph{Written on the Wall II}~\cite{WOWII}; a curated selection is maintained
by West~\cite{WestReg}.  Several consecutive entries concern the maximum
order of an induced tree, an invariant systematically studied by Erd\H{o}s,
Saks, and S\'os~\cite{ESS86}.

Let $G$ be a finite simple connected graph.  The \emph{periphery}
$\Per(G)$ is the set of vertices of eccentricity $\diam(G)$, and the
set-eccentricity of the periphery is
\[
 f(G):=\max_{x\in V(G)} d(x,\Per(G)).
\]
Conjecture~142 of Written on the Wall II asserts
\[
  \tr(G)\ge \frac23\gr(G)+f(G).
\]
Here and in the formal statement, the girth of an acyclic graph is taken to
be zero.  We prove the following stronger cyclic statement.

\begin{theorem}[Graffiti.pc Conjecture 142]\label{thm:main}
Let $G$ be a finite simple connected graph containing a cycle.  If
$g=\gr(G)$ and $f=f(G)$, then
\[
  \tr(G)\ge f+\left\lceil\frac{2g}{3}\right\rceil .
\]
Consequently every finite simple connected graph satisfies
$\tr(G)\ge \frac23\gr(G)+f(G)$.
\end{theorem}

The proof is self-contained apart from elementary facts about finite graphs.
Its structural core is Lemma~\ref{lem:rooted-cycle}, which turns three metric
terminals into an admissible forest attached to a shortest cycle.  The
resulting metric inequality, Lemma~\ref{lem:three-point}, is then combined
with a short classification of the equality case in which the admissible
forest has only one vertex.

\section{Preliminaries}

Write $D=\diam(G)$ and $B=\Per(G)$.  A shortest path is induced, so
\begin{equation}\label{eq:diam-path}
  \tr(G)\ge D+1.
\end{equation}
A shortest cycle $K$ is chordless and isometric: for $u,v\in V(K)$, their
graph distance equals their shorter arc-distance in $K$.  Deleting one
vertex of $K$ therefore gives
\begin{equation}\label{eq:cycle-base}
  \tr(G)\ge g-1.
\end{equation}
If $f\ge1$, then
\begin{equation}\label{eq:D-f}
  D\ge f+1.
\end{equation}
Indeed, always $f\le D$; equality would give a vertex $x$ at distance $D$
from every peripheral vertex, hence $\ecc(x)=D$ and $x\in B$, contradicting
$f=d(x,B)\ge1$.

For a shortest cycle $K$, define $M(K)$ to be the maximum of $|F|$ over all
sets $F\subseteq V(G)\setminus V(K)$ for which $G[F]$ is a forest and there
is a vertex $z\in V(K)$ such that every component $C$ of $G[F]$ sends
exactly one edge into $V(K)\setminus\{z\}$.  Edges from $C$ to $z$ are
unrestricted.

\begin{lemma}[Cycle--forest extension]\label{lem:cycle-forest}
For every shortest cycle $K$,
\[
  \tr(G)\ge g-1+M(K).
\]
\end{lemma}

\begin{proof}
Choose $F,z$ realizing $M(K)$ and let $c$ be the number of components of
$G[F]$.  The graph induced by
$(V(K)\setminus\{z\})\cup F$ is connected.  Its edges consist of the $g-2$
edges of the path $K-z$, the $|F|-c$ forest edges, and one attachment edge
for each of the $c$ components.  Thus it has $g-1+|F|$ vertices and one
fewer edge, and hence is a tree.
\end{proof}

\begin{lemma}[An extra cycle vertex]\label{lem:extra}
Suppose $f\ge1$.  If $g\ge4$, then $\tr(G)\ge g$; if $g\ge5$, then
$M(K)\ge1$ for every shortest cycle $K$.
\end{lemma}

\begin{proof}
If $V(G)=V(K)$, chordlessness and connectedness give $G=C_g$, whose every
vertex is peripheral, contrary to $f\ge1$.  Hence some vertex $y\notin K$
is adjacent to $K$.  When $g\ge5$, it has exactly one neighbour on $K$:
two neighbours and the shorter arc between them would form a cycle of
length at most $\lfloor g/2\rfloor+2<g$.  Taking $F=\{y\}$ and deleting a
cycle vertex other than its root proves both assertions.  For $g=4$, the
only additional possibility is that $y$ has two antipodal neighbours on
$K$; deleting one of them again leaves an induced tree on four vertices.
\end{proof}

We shall also need two elementary connector lemmas.

\begin{lemma}[Three terminals]\label{lem:three-tree}
Every three distinct vertices of a connected triangle-free graph lie in a
common induced tree.
\end{lemma}

\begin{proof}
Let $P$ be an $a$--$b$ geodesic and let $R=r_0\cdots r_\ell$ be a shortest
path from $c=r_0$ to $P$, where $c\notin P$.  Put $q=r_{\ell-1}$ and
$R^\circ=R-r_\ell$.  Only $q$ can have additional neighbours on $P$.
Any two such neighbours have $P$-distance at most two, and triangle-freeness
excludes distance one; hence $q$ has at most two neighbours on $P$.
With one neighbour, $P\cup R^\circ$ induces a tree.  With two neighbours
$u,v$, their $P$-subpath is $u m v$; deleting $m$ joins the two components
of $P-m$ through $q$ and yields a connected induced graph with one fewer
edge than vertices.  Since $m\notin\{a,b\}$, this is the required tree.
\end{proof}

\begin{lemma}[Cycles in a minimal connector]\label{lem:minimal-connector}
Let $J$ be connected and $Q\subseteq V(J)$.  If no proper induced connected
subgraph of $J$ contains $Q$, then every cycle $L$ of $J$ has
$|V(L)|\le |Q|$.
\end{lemma}

\begin{proof}
For every $q\in Q$, choose a shortest path from $q$ to $L$, and denote its
unique first vertex on $L$ by $c(q)$.  Fix $v\in L$.  The path $L-v$ lies
in one component $J_v$ of $J-v$.  Minimality implies that some
$q_v\in Q$ is outside $J_v$.  Its chosen path to $L$ must therefore end at
$v$, so $c(q_v)=v$.  Distinct vertices of $L$ give distinct terminals,
which injects $V(L)$ into $Q$.
\end{proof}

\section{A rooted shortest-cycle lemma}

\begin{lemma}[Rooted shortest cycle]\label{lem:rooted-cycle}
Assume $g\ge5$, and let $S\subseteq V(G)$ with $|S|\le3$.  There exist a
shortest cycle $K$, a vertex $z\in V(K)\setminus S$, and a set
$F\subseteq V(G)\setminus V(K)$ such that
\[
 S\subseteq (V(K)\setminus\{z\})\cup F,
\]
$G[F]$ is a forest, and every component $C$ of $G[F]$ satisfies
\[
 e(C,V(K)\setminus\{z\})=1.
\]
\end{lemma}

\begin{proof}
The empty case is immediate.  Fix an arbitrary shortest cycle $K$, and
choose $X\supseteq V(K)\cup S$ of minimum cardinality such that
$H:=G[X]$ is connected.

Every component $C$ of $H-V(K)$ contains a terminal.  Indeed, it has an
edge to $K$, since $H$ is connected; if $S_C:=S\cap V(C)$ were empty, the
whole component could be deleted while preserving connectedness and all
required vertices.  Put
\[
 A_C:=\{a\in V(C):N_G(a)\cap V(K)\ne\varnothing\}.
\]
Form $J_C$ from $G[C]$ by adjoining a new vertex $r_C$ adjacent precisely
to $A_C$.  This graph is vertex-minimal among induced connected subgraphs
containing $S_C\cup\{r_C\}$.  Otherwise, replacing $C$ by a smaller
connector would leave every resulting component attached to $K$ and would
contradict the choice of $X$.

Lemma~\ref{lem:minimal-connector} shows that every cycle in $J_C$ has length
at most $|S_C|+1\le4$.  A cycle contained in $G[C]$ would also be a cycle of
$G$ and would have length at least $g\ge5$.  Consequently
\begin{equation}\label{eq:C-tree}
  G[C]\text{ is a tree}.
\end{equation}

Let $R_C$ be the minimal subtree of $C$ containing $A_C$.  We claim
\begin{equation}\label{eq:R-small}
  |V(R_C)|\le |S_C|.
\end{equation}
For $v\in V(R_C)$, let $Q_v$ be the component of $J_C-v$ containing $r_C$.
Minimality supplies $s_v\in S_C\setminus V(Q_v)$.  For $s\in S_C$, let
$c(s)$ be the first vertex of $R_C$ on the unique path from $s$ to $R_C$.
Every nonempty component of $R_C-v$ contains a vertex of $A_C$, by the
minimality of $R_C$.  Thus, if $v\ne c(s)$, the path from $s$ to $c(s)$ and
then to such an attachment gives an $s$--$r_C$ path in $J_C-v$; hence
$s\in Q_v$.  Applying this to $s_v$ gives $c(s_v)=v$.  The map
$v\mapsto s_v$ is injective, proving~\eqref{eq:R-small}.

Every $a\in A_C$ has exactly one neighbour on $K$.  Existence follows from
the definition, while two distinct neighbours and their shorter $K$-arc
would form a cycle of length at most $\lfloor g/2\rfloor+2<g$.  Denote the
unique root by $\rho(a)$.  If $a,a'\in A_C$ are distinct, then
\begin{equation}\label{eq:attachment-distance}
 d_C(a,a')\le |V(R_C)|-1\le2
\end{equation}
and their roots are distinct: equal roots, together with the
$a$--$a'$ path in $C$, would form a cycle of length at most four.  Moreover,
the two root edges, that path, and a shorter arc of $K$ give
\begin{equation}\label{eq:girth-roots}
 d_K(\rho(a),\rho(a'))+d_C(a,a')+2\ge g.
\end{equation}

These inequalities sharply restrict the number of attachments.  If
$g\ge9$, then two attachments would make the left side of
\eqref{eq:girth-roots} at most $\lfloor g/2\rfloor+4<g$, so there is only
one.  For $g\in\{7,8\}$, two attachments cannot be adjacent in $C$; three
attachments would fill the three vertices of $R_C$ and contain an adjacent
pair.  For $g=6$, if $R_C=p-q-r$ consisted of three attachments, applying
\eqref{eq:girth-roots} to $p,q$ and to $q,r$ would force both $\rho(p)$ and
$\rho(r)$ to be the unique antipode of $\rho(q)$ on the $6$-cycle,
contrary to distinctness.  Thus every component has at most two
attachments, except that for $g=5$ one component may have three.  Also, at
most one component has two or more attachments, because such a component
contains at least two terminals and $|S|\le3$.

Suppose first that every component has one attachment $a_C$.  The set
\[
 Q=(S\cap V(K))\cup\{\rho(a_C):C\text{ a component of }H-V(K)\}
\]
has cardinality at most $|S|\le3$.  Since $g\ge5$, choose $z\in V(K)\setminus
Q$ and take $F=X\setminus V(K)$.  Equation~\eqref{eq:C-tree} gives a forest,
each component has its single attachment in $Q$, and $z\notin S$.

Next suppose that one component $C_0$ has exactly two attachments with roots
$u,v$.  It contains at least two terminals, so the set
\[
 Q=(S\cap V(K))\cup
   \{\rho(a_C):C\ne C_0\}
\]
has size at most one.  Choose $z\in\{u,v\}\setminus Q$ and again put
$F=X\setminus V(K)$.  Every ordinary component has one edge into $K-z$,
while $C_0$ has exactly the two root edges and loses exactly one of them.
Again $z\notin S$.

It remains to handle three attachments.  Necessarily $g=5$,
$|V(R_C)|=|S_C|=|S|=3$, there are no other off-cycle components, and
$R_C=p-q-r$ with all three vertices attachments.  Write
\[
 u=\rho(p),\qquad v=\rho(q),\qquad w=\rho(r).
\]
Equation~\eqref{eq:girth-roots} gives $d_K(u,v)=d_K(v,w)=2$.  In cyclic
order, write $K=u\alpha v\beta w u$.  Then
\[
  K'=u\alpha v q p u
\]
is another shortest $5$-cycle.

The injection used to prove~\eqref{eq:R-small} is now a bijection.  Hence
there are distinct terminals $s_p,s_q,s_r$ with $c(s_t)=t$.  Let $L_t$ be
the unique $t$--$s_t$ path in the tree $C$, and set
\[
 F'=(V(L_p)\setminus\{p\})\cup
    (V(L_q)\setminus\{q\})\cup V(L_r),
 \qquad z=\alpha.
\]
The three displayed path-pieces are pairwise disjoint and anticomplete, for
an edge between two of them would create a cycle in the tree $C$.  The first
two nonempty pieces attach once to $K'$ at $p$ and $q$, respectively.  The
third attaches once through the edge $rq$; its other old-cycle edge $rw$
does not meet $K'$.  Thus $F'$ is admissible for $(K',z)$.  Finally each
$s_t$ belongs either to its path-piece or, when $s_p=p$ or $s_q=q$, directly
to $K'-z$.  This completes the exceptional case and the proof.
\end{proof}

\section{The three-point bridge}

Let
\[
  \mu:=\max\{M(K):K\text{ is a shortest cycle of }G\}.
\]

\begin{lemma}[Three-point inequality]\label{lem:three-point}
If $g\ge5$, then every three distinct vertices $a,b,c$ satisfy
\begin{equation}\label{eq:three-point}
 d(a,b)+d(b,c)+d(c,a)\le g+2\mu.
\end{equation}
\end{lemma}

\begin{proof}
Apply Lemma~\ref{lem:rooted-cycle} to $S=\{a,b,c\}$ and obtain $K,z,F$.
Retain all edges of $K$, all edges inside the components of $G[F]$, and each
component's unique edge into $K-z$; omit possible edges to $z$.  The resulting
spanning subgraph $U$ of $G[K\cup F]$ is a single cycle $K$ with rooted trees
attached.  It has exactly $|F|$ edges outside $K$: a component with $k$
vertices contributes $k-1$ internal edges and one root edge.

For each pair of terminals, use their unique paths to their roots and a
shortest arc of $K$ between the roots.  An off-cycle edge separates from its
root $q$ of the three terminals.  If its rooted tree contains $h$ terminals,
then its total multiplicity in the three selected paths is
\[
  q(h-q)+q(3-h)=q(3-q)\le2.
\]
Thus the total off-cycle contribution is at most $2|F|$.

The sum of the three pairwise distances between any three roots on a cycle of
length $g$ is at most $g$.  Indeed, if the roots are distinct and the three
intervening arc-lengths are all at most $g/2$, their sum is exactly $g$; if
one arc exceeds $g/2$, the other two provide the shorter route for its
endpoints and the sum is smaller.  Coincident roots are immediate.  Hence
the three chosen paths have total length at most
$g+2|F|\le g+2M(K)\le g+2\mu$.  Graph distances can only be shorter.
\end{proof}

\begin{corollary}\label{cor:metric-mu}
Assume $g\ge5$ and $f\ge1$.  Then
\begin{equation}\label{eq:metric-mu}
  2\mu+g\ge3f+1.
\end{equation}
\end{corollary}

\begin{proof}
Choose $x$ with $d(x,B)=f$ and a diametral pair $b,w$.  Both $b,w$ are
peripheral, so $d(x,b),d(x,w)\ge f$ and $d(b,w)=D\ge f+1$ by
\eqref{eq:D-f}.  Apply Lemma~\ref{lem:three-point} to $x,b,w$.
\end{proof}

The estimate is almost sufficient by itself.  The only delicate integer
case is disposed of by the following sharp classification.

\begin{lemma}[The case $\mu=1$]\label{lem:mu-one}
Assume $g\ge5$ and $f\ge1$.  If $\mu=1$, then $G$ is obtained from $C_g$ by
attaching one pendant vertex to one cycle vertex.  Consequently
\begin{equation}\label{eq:tadpole-f}
  f\le\left\lfloor\frac{g+2}{4}\right\rfloor.
\end{equation}
\end{lemma}

\begin{proof}
Fix a shortest cycle $K$.  Since $M(K)\le\mu=1$, no vertex can have distance
two from $K$: the first two off-cycle vertices on a geodesic to $K$ would
form an admissible two-vertex path.  Thus every vertex outside $K$ is
adjacent to $K$, and by $g\ge5$ it has a unique root.

There cannot be two such vertices $u,v$.  If they are nonadjacent, choose
$z$ different from both roots; the two singleton components are admissible.
If they are adjacent, their roots are distinct (otherwise there is a
triangle), and choosing $z$ to be one root makes $\{u,v\}$ a connected
admissible component.  Either way $M(K)\ge2$, a contradiction.  Since
$f\ge1$, there is exactly one outside vertex, and $G$ is the stated tadpole.

Put $k=\lfloor g/2\rfloor$, let $y$ be the pendant vertex, and let $v_0$ be
its root.  The diameter is $k+1$; the periphery consists of $y$ and the cycle
vertices at cycle-distance $k$ from $v_0$.  Every cycle vertex at position
$i$ on one of the $v_0$--antipode arcs is at distance at most
\[
  \min(i+1,k-i)\le\left\lfloor\frac{k+1}{2}\right\rfloor
  =\left\lfloor\frac{g+2}{4}\right\rfloor
\]
from the displayed peripheral vertices.  The pendant vertex itself is
peripheral, proving~\eqref{eq:tadpole-f}.
\end{proof}

\section{Completion of the proof}

\begin{proposition}\label{prop:large-girth}
Theorem~\ref{thm:main} holds when $g\ge5$ and $f\ge1$.
\end{proposition}

\begin{proof}
Put $a=\lfloor g/3\rfloor$.  Since
$\lceil2g/3\rceil=g-a$, the target is
\begin{equation}\label{eq:target-a}
  \tr(G)\ge f+g-a.
\end{equation}
If $f\le a-1$, the cycle bound~\eqref{eq:cycle-base} gives
$\tr(G)\ge g-1\ge f+g-a$.  Hence write $f=a+s$ with $s\ge0$.

Lemma~\ref{lem:extra} gives $\mu\ge1$, while
Lemma~\ref{lem:cycle-forest} gives $\tr(G)\ge g-1+\mu$.  It therefore
suffices to prove
\begin{equation}\label{eq:mu-goal}
  \mu\ge s+1.
\end{equation}
Corollary~\ref{cor:metric-mu} and the three residues of $g$ give:
\begin{align*}
 g=3a   &: &2\mu&\ge3s+1,\\
 g=3a+1 &: &2\mu&\ge3s,\\
 g=3a+2 &: &2\mu&\ge3s-1.
\end{align*}
In the first case, $\mu\le s$ would imply $2\mu\le2s<3s+1$.
In the second, $s=0$ is covered by $\mu\ge1$, while for $s\ge1$ we have
$2s<3s$.  In the third, $s=0$ is again immediate, and $s\ge2$ gives
$2s<3s-1$.

The only remaining possibility is $g=3a+2$ and $s=1$.  If $\mu=1$,
Lemma~\ref{lem:mu-one} yields
\[
 f\le\left\lfloor\frac{3a+4}{4}\right\rfloor\le a,
\]
because $a\ge1$, contradicting $f=a+1$.  Hence $\mu\ge2=s+1$.
This proves~\eqref{eq:mu-goal}, and therefore~\eqref{eq:target-a}.
\end{proof}

\begin{proposition}\label{prop:small-girth}
Theorem~\ref{thm:main} holds for $g=3$ and $g=4$ when $f\ge1$.
\end{proposition}

\begin{proof}
For $g=3$, equations~\eqref{eq:diam-path} and~\eqref{eq:D-f} give
$\tr(G)\ge D+1\ge f+2=f+\lceil2g/3\rceil$.

Let $g=4$.  Choose an $f$-realizer $x$ and a diametral pair $b,w$.  The graph
is triangle-free, so Lemma~\ref{lem:three-tree} supplies an induced tree $T$
containing $x,b,w$.  Let $T_0$ be the minimal subtree containing them.  Every
edge of $T_0$ lies on exactly two of the three terminal-to-terminal paths;
therefore
\[
  2|E(T_0)|=d_T(x,b)+d_T(x,w)+d_T(b,w)
   \ge 2f+D\ge3f+1.
\]
If $f\ge2$, then
\[
 |V(T)|\ge1+\left\lceil\frac{3f+1}{2}\right\rceil\ge f+3.
\]
If $f=1$, Lemma~\ref{lem:extra} gives $\tr(G)\ge g=4=f+3$.
\end{proof}

\begin{proof}[Proof of Theorem~\ref{thm:main}]
The preceding two propositions cover cyclic graphs with $f\ge1$.  If $G$ is
cyclic and $f=0$, then $g-1\ge\lceil2g/3\rceil$ for $g\ge3$, so
\eqref{eq:cycle-base} applies.  This proves the integral cyclic assertion.

If $G$ is acyclic, connectedness makes $G$ itself a tree, so
$\tr(G)=|V(G)|$ and $f\le D\le |V(G)|-1$ (with the one-vertex case
immediate).  Since $\gr(G)=0$, the real-valued Formal Conjectures inequality
also follows.  Finally,
$\lceil2g/3\rceil\ge2g/3$, completing the proof.
\end{proof}

\section{Verification and formalization status}\label{sec:verification}

The proof above was checked in two independent forms.  First, a separate
constructive proof based on descents to an arbitrary shortest cycle was
implemented literally: every branch produces its cycle vertex $z$, its
admissible forest, or its induced path, and a fresh checker verifies
inducedness, connectedness, acyclicity, edge multiplicities, and the claimed
cardinality.  On a corpus of 10,776 connected cyclic graphs (the complete
NetworkX atlas through seven vertices, structured extremal families, and
seeded random and adversarial graphs through 32 vertices), the validator
reported no failure.  Repeating every free choice with two independent
random seeds again gave no failure, for 32,328 checked certificates in all.
These computations are supporting checks, not a substitute for the proof.

The exact statement of Conjecture~142 is formalized in the Google DeepMind
\emph{Formal Conjectures} repository~\cite{FC,FC142}, using the repository's
largest-induced-tree invariant, Mathlib's natural-valued girth, and the
set-eccentricity of the boundary (peripheral) vertices.  A compiled Lean~4
skeleton already closes the acyclic and $f=0$ branches, and the reusable
cycle--forest extension API has been checked.  At the date of this draft, the
rooted shortest-cycle lemma and the final cyclic branch have not yet been
fully formalized; we therefore make no claim here of a completed Lean proof.

\begin{remark}[Sharpness]
Equality occurs throughout the $g=3$ diameter branch.  Among the exhaustively
checked small graphs, the only equality class of girth greater than three is
the $6$-cycle with one pendant vertex: here $f=2$ and
$\tr(G)=6=f+\lceil2g/3\rceil$.  In the proof this is exactly the
$\mu=1$ tadpole case of Lemma~\ref{lem:mu-one}.
\end{remark}

\paragraph{Acknowledgements.}
The author thanks Ermelinda DeLaVi\~na and Douglas B.\ West for maintaining
the Graffiti.pc conjecture lists, and the maintainers of the Formal
Conjectures repository for the formal statement.  OpenAI GPT Pro and Codex
were used to explore proof routes, run finite checks, and assist with editing;
the mathematical argument presented here was separately subjected to
adversarial proof review.

\begin{thebibliography}{10}

\bibitem{DeL01}
E.~DeLaVi\~na,
\emph{Graffiti.pc},
Graph Theory Notes of New York XLII:3 (2002), 26--30.

\bibitem{DeLhistory}
E.~DeLaVi\~na,
\emph{Some history of the development of Graffiti},
in: Graphs and Discovery, DIMACS Ser.\ Discrete Math.\ Theoret.\ Comput.\
Sci.\ 69, AMS, 2005, 81--118.

\bibitem{WOWII}
E.~DeLaVi\~na,
\emph{Written on the Wall II: Conjectures of Graffiti.pc},
web list, University of Houston--Downtown,
\url{http://cms.dt.uh.edu/faculty/delavinae/research/wowII/}
(retrieved July 2026).

\bibitem{WestReg}
D.~B. West,
\emph{Some conjectures of Graffiti.pc}, web register,
\url{https://dwest.web.illinois.edu/regs/graffiti.html}
(retrieved July 2026).

\bibitem{Faj88}
S.~Fajtlowicz,
\emph{On conjectures of Graffiti},
Discrete Math.\ 72 (1988), 113--118.

\bibitem{ESS86}
P.~Erd\H{o}s, M.~Saks, and V.~T. S\'os,
\emph{Maximum induced trees in graphs},
J.~Combin.\ Theory Ser.\ B 41 (1986), 61--79.
\href{https://doi.org/10.1016/0095-8956(86)90028-6}
{doi:10.1016/0095-8956(86)90028-6}.

\bibitem{DW04}
E.~DeLaVi\~na and B.~Waller,
\emph{On some conjectures of Graffiti.pc on the maximum order of induced
subgraphs},
Congr.\ Numer.\ 166 (2004), 11--32.

\bibitem{FC}
Google DeepMind,
\emph{Formal Conjectures}, GitHub repository,
\href{https://github.com/google-deepmind/formal-conjectures}{repository homepage}.

\bibitem{FC142}
Formal Conjectures,
\emph{GraphConjecture142.lean},
\href{https://github.com/google-deepmind/formal-conjectures/blob/main/FormalConjectures/WrittenOnTheWallII/GraphConjecture142.lean}{source file on GitHub}
(retrieved July 2026).

\bibitem{Lean4}
L.~de~Moura and S.~Ullrich,
\emph{The Lean 4 theorem prover and programming language},
in: Automated Deduction -- CADE 28, LNCS 12699, Springer, 2021, 625--635.

\end{thebibliography}

\end{document}
